\extrachap{Predgovor}

S povečevanjem gostote integriranih vezij postajajo tranzistorji vse manjše nanoelektronske naprave. Želeno delovanje tranzistorjev velikosti le nekaj nanometrov pa pričnejo ovirati kvantni pojavi in sile med elektroni. Zato so se pojavile zamisli o novih procesnih platformah, ki jih ti pojavi ne ovirajo, temveč jih izkoriščajo za procesiranje. Nove nanoelektronske naprave lahko prinesejo tudi večjo gostoto, manjšo porabo energije in hitrejše delovanje.

Ena od takih platform so kvantni celični avtomati (angl.\ \emph{quantum-dot cellular automata}, QCA)~\cite{lent1993, lent1997}. Namesto tranzistorjev vsebujejo celice z elektroni, katerih razporeditev znotraj celice določa logično stanje. Procesiranje se izvede z ustrezno razporeditvijo celic, pri čemer se izkoriščata Coulombova sila in kvantno tuneliranje. Fizično so bile doslej izdelane le manjše strukture, od prvih majoritetnih vrat iz kovinskih kvantnih pik~\cite{amlani1999} do vodil in vrat iz posameznih atomov na površini silicija~\cite{huff2018}. Večja vezja lahko načrtujemo, simuliramo in optimiziramo z računalniškimi orodji~\cite{walus2004, silltorres2018, walter2019fiction}.

Kvantno procesiranje z osnovnimi gradniki, kvantnimi biti, temelji na superpoziciji, interferenci in prepletenosti, ki obljubljajo eksponentno pohitritev nekaterih algoritmov glede na klasične računalnike~\cite{nielsen2010}. V zadnjih letih so prikazali hitrejše reševanje posebej izbranih problemov~\cite{arute2019quantum, wu2021strong}. Današnji kvantni procesorji so še vedno občutljivi na šum, vendar se izboljšuje tudi popravljanje napak~\cite{google2025below}. Simulatorji in dostop do manjših kvantnih računalnikov prek spleta omogočajo lažje preučevanje področja.

Kvantnim celičnim avtomatom in kvantnemu procesiranju je skupna reverzibilnost. Fizično reverzibilno procesiranje brez izgube informacij se približa najmanjši možni porabi energije~\cite{landauer1961, bennett1973}. Kvantno procesiranje je inherentno reverzibilno, saj so kvantna vrata unitarna, ustrezna zasnova vezja QCA pa omogoča fizično reverzibilno procesno platformo.

Zbornik sestavljajo poglavja, ki jih v okviru seminarskih nalog pripravijo študenti. Vsako poglavje odgovarja na lastno raziskovalno vprašanje in temelji na poskusih, ki so jih avtorji izvedli sami.
